% -------------------------------------------------------------------------
% WBS ID: RES-GEO-TAX
% Deliverable: Barrier Geometry Taxonomy and Candidate Families
% Status: In progress
% Evidence status: Barrier-class scope established; real-defence source gap open
% Review status: Not reviewed
% Last updated: 2026-07-19
% Dependencies: RES-DAT-SCN, RES-DAT-CAS, RES-PHY-BAR
% Downstream interfaces: RES-GEO-PAR, RES-GEO-MET, RES-GEO-REP, Software geometry manifest
% -------------------------------------------------------------------------

\section{RES-GEO-TAX --- Barrier Geometry Taxonomy and Candidate Families}
\label{sec:res-geo-tax}

\subsection{Purpose and Scope}
\label{sec:res-geo-tax-purpose}

This Deliverable classifies the active fixed-rigid barrier families
used for the Tohoku corridor comparison. It admits two baseline
families:

\begin{enumerate}
  \item wall-type barriers, including continuous or segmented seawalls,
    ridges, breakwaters and raised impermeable lines;
  \item obstacle, dissipating or array-type barriers, including finite
    blocks, columns, staggered rows, permeable-looking resolved
    openings and other structures whose effect is controlled by
    blockage, spacing and wake interaction.
\end{enumerate}

Selection is based on geometric recoverability, hydrodynamic relevance
and numerical representability. The deliverable excludes unrestricted
geometry generation and does not treat damage, scour or deformable
barrier response as active baseline geometry.

\subsection{Research Questions}
\label{sec:res-geo-tax-research-questions}

\begin{enumerate}
  \item Which barrier families are active in the Kamaishi and Sendai
    comparison cases?
  \item Which minimum descriptors distinguish wall-type and
    obstacle/array-type barriers for the local URANS--VOF model?
  \item Which real-defence geometries have enough provenance to be
    reconstructed rather than treated as TODO candidates?
\end{enumerate}

\subsection{Terminology and Definitions}
\label{sec:res-geo-tax-definitions}

\begin{description}
  \item[Wall-type barrier] A laterally extended defence whose principal
    hydrodynamic effect is reflection, run-up, overtopping and
    blocking along a near-continuous crest.
  \item[Obstacle or dissipating barrier] A finite-width body or row of
    bodies whose principal hydrodynamic effect is local blockage,
    separation, wake production, transmission and turbulent energy
    loss.
  \item[Array] A repeated set of resolved obstacle elements with
    spacing, staggering and row count recorded explicitly.
  \item[Fixed rigid barrier] A solid geometry that supplies no-slip and
    no-penetration wall patches to the local fluid solver and does not
    deform during the baseline hydrodynamic calculation.
\end{description}

\subsection{Literature Evidence}
\label{sec:res-geo-tax-literature-evidence}

\textcite{deljesus2012porous} supports the taxonomy split between
impermeable vertical barriers and barriers whose porosity or openings
require explicit modelling assumptions. Source role: supports the
wall-type versus dissipating/opening class distinction and limits any
claim about unresolved porous resistance. Project conclusion:
porosity-like behaviour is active only when openings are resolved by
geometry or a later porous closure is formally activated.

\textcite{watanabe2022validation} supports using controlled idealised
coastal structures for local hydrodynamic validation. Source role:
supports simple fixed obstacle and wall geometries as validation-facing
families rather than arbitrary generated forms. Limitation: the paper
does not define the final Kamaishi or Sendai defence geometry.

\textcite{mori2012nationwide} supports retaining Sanriku/Kamaishi and
Sendai coastal-plain settings as contrasting Tohoku impact regimes.
Source role: motivates site-facing geometry classes, not detailed
barrier dimensions. TODO/missing-source: identify a citable primary
source for the Kamaishi bay-mouth breakwater geometry, 2011 damage
state and any reconstruction dimensions before recreating it as a
project barrier.

\subsection{Governing Relations and Quantitative Definitions}
\label{sec:res-geo-tax-governing-relations}

For software-facing comparison, each geometry record shall include a
class label and a minimum descriptor vector:

\begin{align}
  \mathcal{G}_j
  &=
  \left\{
    c_j,\,
    \mathbf{x}_j,\,
    \psi_j,\,
    H_j,\,
    B_j,\,
    L_j,\,
    C_j,\,
    S_j,\,
    N_j,\,
    \chi_j
  \right\}
  \label{eq:res-geo-tax-minimum-geometry-record}
\end{align}

where $c_j$ is the barrier class, $\mathbf{x}_j$ is the placement
reference point in the local corridor coordinate system, $\psi_j$ is
the plan-view orientation, $H_j$ is height, $B_j$ is width or
thickness, $L_j$ is along-barrier length, $C_j$ is crest elevation,
$S_j$ is spacing for arrays, $N_j$ is element or row count and
$\chi_j$ records resolved openings, roughness and provenance flags.
Decision role: \cref{eq:res-geo-tax-minimum-geometry-record} is a data
contract, not a design-generation equation. Supporting evidence:
\textcite{deljesus2012porous} and \textcite{watanabe2022validation}.
Limitation: final admissible parameter bounds remain site- and
mesh-dependent.

\subsection{Geometry Family or Representation}
\label{sec:res-geo-tax-geometry-family-representation}

The active family set is:

\begin{itemize}
  \item \texttt{wall}: continuous or segmented impermeable line
    defence, with optional resolved openings;
  \item \texttt{obstacle}: isolated block, pier, mound or short
    barrier element;
  \item \texttt{array}: repeated obstacle elements whose spacing,
    staggering and row count are explicit;
  \item \texttt{dissipating}: resolved obstacle or array geometry
    intended to increase separation, wake loss, reflection or
    overtopping reduction.
\end{itemize}

Unresolved porous-media barriers, compliant structures, scour-adaptive
foundations and damage-evolving barriers are deferred extensions.

\subsection{Design Variables and Parameter Bounds}
\label{sec:res-geo-tax-design-variables-parameter-bounds}

% TODO: Complete this RES-GEO Domain-specific subsection.

\subsection{Geometric Descriptors}
\label{sec:res-geo-tax-geometric-descriptors}

% TODO: Complete this RES-GEO Domain-specific subsection.

\subsection{Placement and Arrangement Rules}
\label{sec:res-geo-tax-placement-arrangement-rules}

% TODO: Complete this RES-GEO Domain-specific subsection.

\subsection{Feasibility Constraints}
\label{sec:res-geo-tax-feasibility-constraints}

% TODO: Complete this RES-GEO Domain-specific subsection.

\subsection{Two- and Three-Dimensional Representation}
\label{sec:res-geo-tax-two-three-dimensional-representation}

% TODO: Complete this RES-GEO Domain-specific subsection.

\subsection{Candidate Approaches}
\label{sec:res-geo-tax-candidate-approaches}

% TODO: Define the credible candidate approaches considered.

\subsection{Critical Comparison}
\label{sec:res-geo-tax-critical-comparison}

% TODO: Compare candidates using physically and project-relevant criteria.
% Do not select an approach solely on popularity or implementation ease.

\subsection{Assumptions and Applicability}
\label{sec:res-geo-tax-assumptions}

% TODO: State assumptions and the conditions under which they are valid.

\subsection{Limitations and Failure Conditions}
\label{sec:res-geo-tax-limitations}

% TODO: State where the evidence, model or selected approach ceases to be
% reliable or applicable.

\subsection{Project-Specific Conclusions}
\label{sec:res-geo-tax-conclusions}

The project taxonomy is now class-comparison led. A geometry may enter
the active campaign when it has a traceable descriptor record, a fixed
rigid local representation and a clear hydrodynamic purpose. Real
defence recreation is allowed only after source geometry is identified;
otherwise the case remains a TODO rather than an invented citation.

\subsection{Selected Approach and Rationale}
\label{sec:res-geo-tax-selected-approach}

The selected approach is a small active taxonomy: wall-type barriers
and obstacle/dissipating/array barriers. The rejected active approach
is unrestricted generated geometry, because the present work must first
stabilise data handling, local hydrodynamic load extraction and the
V1--V2--V3 validation hierarchy.

\subsection{Unresolved Decisions}
\label{sec:res-geo-tax-unresolved-decisions}

\begin{itemize}
  \item TODO: obtain a primary or peer-reviewed source for Kamaishi
    bay-mouth breakwater dimensions and 2011 condition.
  \item Select final descriptor bounds after RES-DAT-GEO terrain
    extraction and RES-NUM-MSH mesh constraints are available.
  \item Decide whether any dissipating class requires a porous closure
    rather than resolved openings.
\end{itemize}

\subsection{Requirements and Handoffs}
\label{sec:res-geo-tax-requirements}

Software shall support a geometry manifest containing the descriptor
record in \cref{eq:res-geo-tax-minimum-geometry-record}, class labels,
local corridor coordinates, source/provenance fields and flags for
resolved openings versus deferred porous physics. `RES-MOD-IBC` and
`RES-NUM-MSH` shall consume this record as fixed terrain/barrier wall
patches; `RES-VER-MET` shall use the same identifiers for load and
energy metrics.

\subsection{Deliverable Completion Record}
\label{sec:res-geo-tax-completion-record}

% TODO: Record whether the Deliverable is:
%
% - Not started
% - In progress
% - Draft complete
% - Reviewed
% - Accepted
%
% Acceptance requires:
%
% - the research questions to be answered;
% - claims to be supported by appropriate evidence;
% - assumptions and limitations to be explicit;
% - the project-specific conclusion to be stated;
% - downstream requirements to be traceable;
% - unresolved decisions to be closed or formally deferred.
